\documentclass[11pt,a4paper]{article}

% ============================================================
%  Informe breve: entrenamientos SSL 3D en RunPod, split 361/63/84
%  Figuras: python scripts/build_report_split361.py
%  Compilar desde docs/: pdflatex informe_split361.tex (dos veces)
% ============================================================

\usepackage[a4paper,margin=2.2cm]{geometry}
\usepackage[T1]{fontenc}
\usepackage{lmodern}
\usepackage[spanish,es-noshorthands]{babel}
\usepackage{microtype}
\usepackage{xcolor}
\usepackage{amsmath,amssymb}
\usepackage{booktabs}
\usepackage{array}
\usepackage{tabularx}
\usepackage{graphicx}
\usepackage{caption}
\usepackage{float}
\usepackage{placeins}
\usepackage{enumitem}
\usepackage{fancyhdr}
\usepackage{titlesec}
\usepackage[most]{tcolorbox}
\usepackage{tikz}
\usetikzlibrary{arrows.meta,positioning,fit,backgrounds,calc}

% Paleta del informe v2 (VICReg azul, BYOL naranja, DINOv2 aqua)
\definecolor{primaryblue}{HTML}{1B4F72}
\definecolor{steelblue}{HTML}{2E86AB}
\definecolor{accentorange}{HTML}{E67E22}
\definecolor{softgray}{HTML}{85929E}
\definecolor{panelbg}{HTML}{EAF2F8}
\definecolor{warnbg}{HTML}{FDEBD0}
\definecolor{tldrbg}{HTML}{EBF5FB}
\definecolor{vicregc}{HTML}{2A78D6}
\definecolor{byolc}{HTML}{EB6834}
\definecolor{dinoc}{HTML}{1BAF7A}
\definecolor{diagbg}{HTML}{FEF5E7}

\usepackage[colorlinks=true,linkcolor=primaryblue,citecolor=primaryblue,urlcolor=steelblue]{hyperref}

% Cajas
\newtcolorbox{tldrbox}{colback=tldrbg,colframe=primaryblue,coltitle=white,
  colbacktitle=primaryblue,fonttitle=\bfseries,boxrule=0.9pt,arc=3pt,
  left=8pt,right=8pt,top=5pt,bottom=5pt,title=En pocas palabras}
\newtcolorbox{ideabox}[1][Idea clave]{colback=panelbg,colframe=steelblue,coltitle=white,
  colbacktitle=steelblue,fonttitle=\bfseries,boxrule=0.8pt,arc=2pt,
  left=8pt,right=8pt,top=5pt,bottom=5pt,title=#1}
\newtcolorbox{warnbox}[1][Ojo con esto]{colback=warnbg,colframe=accentorange,coltitle=white,
  colbacktitle=accentorange,fonttitle=\bfseries,boxrule=0.8pt,arc=2pt,
  left=8pt,right=8pt,top=5pt,bottom=5pt,title=#1}
\newtcolorbox{byolbox}[1][BYOL]{colback=diagbg,colframe=byolc,
  coltitle=white,colbacktitle=byolc,fonttitle=\bfseries,boxrule=0.8pt,arc=2pt,
  left=8pt,right=8pt,top=5pt,bottom=5pt,title=#1}

% Encabezados y secciones
\pagestyle{fancy}
\fancyhf{}
\renewcommand{\headrulewidth}{0.5pt}
\fancyhead[L]{\small\color{softgray}SSL 3D en RunPod · split 361/63/84}
\fancyhead[R]{\small\color{softgray}VICReg / BYOL / DINOv2}
\fancyfoot[C]{\small\thepage}
\setlength{\headheight}{14pt}

\titleformat{\section}{\normalfont\Large\bfseries\color{primaryblue}}{\thesection}{0.6em}{}
\titleformat{\subsection}{\normalfont\large\bfseries\color{steelblue}}{\thesubsection}{0.6em}{}

\graphicspath{{figures_split361/}}
\setlength{\parskip}{4pt}
\setlength{\parindent}{0pt}
\setlength{\emergencystretch}{3em}
\setlist[itemize]{topsep=2pt,itemsep=2pt}
\setlist[enumerate]{topsep=2pt,itemsep=2pt}
\captionsetup{font=small,labelfont={bf,color=primaryblue}}
\renewcommand{\arraystretch}{1.15}

% Estilos TikZ
\tikzset{
  viewbox/.style={rectangle, rounded corners=1pt, draw=softgray, thick, fill=white,
    text width=2.6cm, align=center, minimum height=0.8cm, inner sep=3pt, font=\scriptsize},
  netbox/.style={rectangle, rounded corners=2pt, draw=#1, thick, fill=#1!8,
    text width=3.4cm, align=center, minimum height=0.85cm, inner sep=3pt, font=\scriptsize},
  targetbox/.style={rectangle, rounded corners=2pt, draw=softgray, thick, dashed,
    fill=softgray!10, text width=3.4cm, align=center, minimum height=0.85cm, inner sep=3pt,
    font=\scriptsize},
  lossbox/.style={rectangle, rounded corners=2pt, draw=accentorange, thick,
    fill=accentorange!8, text width=9.6cm, align=left, inner sep=5pt, font=\scriptsize},
  evalbox/.style={rectangle, rounded corners=2pt, draw=primaryblue, thick, fill=panelbg,
    text width=3.2cm, align=center, minimum height=1cm, inner sep=4pt, font=\scriptsize},
  arr/.style={-{Latex[length=2mm]}, thick, primaryblue},
  arrg/.style={-{Latex[length=2mm]}, thin, softgray},
  emaarr/.style={-{Latex[length=2mm]}, thin, dashed, accentorange},
}

\begin{document}

% ============================================================
% PORTADA
% ============================================================
\begin{titlepage}
\centering
\vspace*{0.8cm}
{\Huge\bfseries\color{primaryblue} SSL 3D en OASIS3}\\[8pt]
{\Large\color{steelblue} Los entrenamientos en RunPod con el split 361/63/84\\
versión corta y visual}\\[20pt]

\begin{tikzpicture}
\draw[primaryblue, thick] (0,0) -- (11,0);
\end{tikzpicture}

\vspace{14pt}
{\large \textit{Qué datos entran, qué ve cada modelo, cómo aprende y cómo lo medimos,\\
en una lectura de diez minutos.}}

\vspace{24pt}
\includegraphics[width=0.95\textwidth]{roi_cubes.png}

\vspace{16pt}
\begin{tcolorbox}[colback=panelbg,colframe=primaryblue,arc=4pt,width=0.9\textwidth,
  boxrule=1pt, left=14pt,right=14pt,top=8pt,bottom=8pt]
\begin{center}
\textbf{oasis3-ssl-workbench · \texttt{campaign\_ssl3d\_runpod/split361\_30}}\\[4pt]
\small 508 sujetos, split por sujeto 361/63/84 · VICReg, BYOL y DINOv2 con su receta campeona\\
30 épocas · 3 semillas · probe lineal sobre el encoder congelado\\
RTX 5090 en RunPod · 9 corridas en 20 minutos · 0.33 USD
\end{center}
\end{tcolorbox}

\vfill
{\small Resultados: \texttt{split361\_30/REPORT.md} y \texttt{DELTAS.md}; diagnóstico de BYOL:
\texttt{byol\_diag30/DIAGNOSTICO.md}.}\\[4pt]
{\small 1 de octubre de 2026}
\end{titlepage}

\tableofcontents
\newpage

% ============================================================
\section{En pocas palabras}
% ============================================================

\begin{tldrbox}
\begin{itemize}[leftmargin=12pt]
\item \textbf{Los tres métodos aprenden y ninguno colapsa.} En sujetos no vistos, la
  diferencia entre sujetos supera entre 10 y 34 veces la variación que introduce la
  augmentation.
\item \textbf{En test los tres empatan}: AUROC de 0.831 a 0.837 sobre 84 sujetos.
\item \textbf{DINOv2 es el que más gana con el SSL}: $+0.055$ en CV y $+0.104$ en test
  frente al mismo ViT sin entrenar, ambos significativos. Pero parte del punto más bajo.
\item \textbf{VICReg tiene el mejor CV} (0.875) y mejora en CV ($+0.026$, significativo).
\item \textbf{BYOL no supera a su inicio ImageNet en CV} ($-0.009$). Su pérdida de
  validación sube, pero no por memorizar: es un efecto de la batch norm (Sección~\ref{sec:byol}).
\item \textbf{Con 100 épocas nada mejora} (Sección~\ref{sec:e100}). El SSL empata en test con el
  mismo encoder entrenado de forma supervisada; frente al encoder sin SSL gana DINOv2 (CV y
  test), VICReg solo en CV y BYOL no.
\end{itemize}
\end{tldrbox}

\begin{table}[H]
\centering
\small
\begin{tabular}{lllrrrr}
\toprule
Método & Entrada & Augmentation & AUROC CV & AUROC test & $\Delta$ CV & $\Delta$ test \\
\midrule
\textcolor{vicregc}{\textbf{VICReg}} & hipocampos 2$\times$32$^3$ & \texttt{geo} & \textbf{0.875} & 0.835 & $+0.026^{*}$ & $+0.029$ \\
\textcolor{byolc}{\textbf{BYOL}} & hipocampos 2$\times$32$^3$ & \texttt{noise} & 0.852 & \textbf{0.837} & $-0.009$ & $+0.029$ \\
\textcolor{dinoc}{\textbf{DINOv2}} & 3 regiones 6$\times$32$^3$ & \texttt{elastic} & 0.839 & 0.831 & $+0.055^{*}$ & $+0.104^{*}$ \\
\bottomrule
\end{tabular}
\caption{Resumen (media de 3 semillas). $\Delta$: frente al mismo encoder y la misma entrada
sin entrenamiento SSL; $^{*}$ = intervalo de confianza del 95\,\% enteramente sobre cero.}
\label{tab:summary}
\end{table}

% ============================================================
\section{Los datos y el split}
\label{sec:data}
% ============================================================

Cada sesión de OASIS3 es una resonancia T1 registrada al espacio MNI con vóxeles de 2\,mm.
El split reparte \textbf{508 sujetos por sujeto}: ninguno aparece en dos conjuntos. Val y test
tienen cupos fijos de HC y AD, y dentro de cada diagnóstico se reparten de forma proporcional
entre los tres terciles de edad (semilla 42). Los terciles se calculan por separado para HC
(cortes en 65 y 73 años) y para AD (72.3 y 79), con la edad de la sesión evaluada. El sexo no
está registrado en el CSV clínico de OASIS3, y los años de estudio no fueron un estrato.

\begin{figure}[H]
\centering
\includegraphics[width=0.9\textwidth]{split_balance.png}
\caption{Sujetos por split. Arriba, diagnóstico: ``otros'' (sin diagnóstico, Lewy,
frontotemporal, vascular o etiqueta ambigua) solo entran al entrenamiento SSL, que no usa
etiquetas. Al medio, los estratos de edad con que se armó el split. Abajo, años de estudio, que no
se estratificaron: val queda algo más alto (15.9) y test algo más bajo (14.8) que train (15.3).}
\label{fig:split}
\end{figure}

\begin{itemize}
\item \textbf{Train (361):} el SSL toma \emph{una sesión al azar por sujeto} en cada época:
  unos 23 lotes de 16, unos 690 pasos en 30 épocas.
\item \textbf{Val (63) y test (84):} solo AD y HC, una sesión cada uno. Val sirve para la
  pérdida de validación y para elegir la regularización del probe; test se usa una sola vez.
\end{itemize}

\begin{warnbox}[La edad viene pegada al diagnóstico]
En los tres splits los AD tienen unos 7 años más que los HC (HC 67--69 años, AD 75--76).
Parte de lo que separa el probe puede ser envejecimiento y no Alzheimer.
\end{warnbox}

% ============================================================
\section{Qué cubo ve cada modelo}
\label{sec:roi}
% ============================================================

Nadie ve el cerebro entero: cada modelo ve \textbf{cubos de 32 vóxeles} (64\,mm) recortados
en posiciones fijas, el centroide poblacional de cada región. Se guarda un cubo de 38 para
tener margen en recortes y deformaciones. El hemisferio derecho se \emph{espeja}: ambos lados
quedan orientados igual y voltear la cabeza equivale a intercambiar canales.

\begin{table}[H]
\centering
\small
\begin{tabularx}{\textwidth}{l>{\raggedright\arraybackslash}Xc>{\raggedright\arraybackslash}X}
\toprule
Modelo & Regiones & Entrada & Por qué \\
\midrule
VICReg y BYOL & hipocampo izq. y der. & 2$\times$32$^3$ & 32$^3$ fue la mejor ROI para las CNN \\
DINOv2 & hipocampo, ventrículo lateral, corteza posteromedial & 6$\times$32$^3$ & 3 regiones frente a 32$^3$: $+0.034$ CV, $+0.077$ test \\
\bottomrule
\end{tabularx}
\caption{ROI por modelo, elegida en la ronda de ROI (50 épocas, split anterior). Para las CNN,
cubos de 48$^3$ rindieron igual y los de 64$^3$ o las 3 regiones, peor.}
\end{table}

\begin{figure}[H]
\centering
\includegraphics[width=\textwidth]{inputs.png}
\caption{Lo que entra a la red, en el corte coronal central de cada canal. La corteza
posteromedial es el cíngulo posterior e ístmico más el precúneo.}
\label{fig:inputs}
\end{figure}

% ============================================================
\section{Augmentations: una receta por modelo}
\label{sec:aug}
% ============================================================

Los métodos SSL aprenden comparando dos \emph{vistas} distintas del mismo sujeto. Cada modelo
usa la receta que ganó en el \emph{screen} de 19 recetas (20 épocas): una sola idea por receta,
más el volteo izquierda-derecha (intercambio de canales, $p=0.5$) que llevan todas.

\begin{ideabox}[Qué cambió frente al informe v2]
Ya no hay rotación con \textbf{relleno de ceros}: sus esquinas negras eran el atajo sobre el
que colapsaba BYOL. Todas las rotaciones rellenan con el borde.
\end{ideabox}

\begin{figure}[H]
\centering
\includegraphics[width=0.84\textwidth]{views.png}
\caption{Cubo limpio que ve el probe y dos vistas de cada receta (canal del hipocampo izquierdo,
sujeto de validación). Cada imagen apila los cortes centrales de los tres planos del mismo
volumen (coronal al frente, axial y sagital detrás): la red no ve cortes sueltos sino el volumen
3D completo, y las augmentations también actúan en 3D.}
\label{fig:views}
\end{figure}

\begin{table}[H]
\centering
\small
\begin{tabularx}{\textwidth}{l>{\raggedright\arraybackslash}X>{\raggedright\arraybackslash}X>{\raggedright\arraybackslash}X}
\toprule
 & \textcolor{vicregc}{\textbf{VICReg · \texttt{geo}}} & \textcolor{byolc}{\textbf{BYOL · \texttt{noise}}} & \textcolor{dinoc}{\textbf{DINOv2 · \texttt{elastic}}} \\
\midrule
Recorte & sub-cubo de 32--35 vóxeles en posición al azar, reescalado a 32 & cubo central de 32 & cubo central de 32 \\
Rotación & $\pm5^\circ$ por eje, relleno con el borde & no & no \\
Intensidad & no & ruido gaussiano, $\sigma\sim U(0;\,0.22)$ & no \\
Deformación & no & no & elástica TorchIO: 7 puntos de control, $\leq$1.95 vóxeles \\
Volteo izq.-der. & $p=0.5$ & $p=0.5$ & $p=0.5$ \\
Enmascarado & no & no & iBOT: 50\,\% de las vistas, 10--50\,\% de los parches \\
\bottomrule
\end{tabularx}
\caption{Recetas (\texttt{t5\_ssl.recipes.AUG\_RECIPES}). La intensidad del ruido y de la
deformación está calibrada para alterar el cubo lo mismo (RMSE 0.08).}
\label{tab:aug}
\end{table}

% ============================================================
\section{Arquitecturas}
\label{sec:arch}
% ============================================================

Las tres redes parten de pesos de \textbf{ImageNet en 2D} y se \emph{inflan} a 3D: cada filtro
$k\times k$ se repite $k$ veces en profundidad y se divide por $k$. En el ViT, el parche de
16$\times$16 se reduce a 8$^3$ y las posiciones se extienden a una grilla de 4$\times$4$\times$4.

\begin{table}[H]
\centering
\small
\begin{tabularx}{\textwidth}{l>{\raggedright\arraybackslash}p{3.0cm}>{\raggedright\arraybackslash}X>{\raggedright\arraybackslash}p{2.1cm}}
\toprule
Método & Encoder & Cabezas (solo para entrenar) & Lo que usa el probe \\
\midrule
\textcolor{vicregc}{\textbf{VICReg}} & ConvNeXt-T 3D & proyector 768$\rightarrow$2048$\rightarrow$2048$\rightarrow$2048 (BN, dropout 0.1) & 768 (GAP) \\
\textcolor{byolc}{\textbf{BYOL}} & ResNet-18 3D sin max-pool & proyector 512$\rightarrow$4096$\rightarrow$256 (BN, dropout 0.1); predictor 256$\rightarrow$4096$\rightarrow$256 (BN) & 512 (GAP, red online) \\
\textcolor{dinoc}{\textbf{DINOv2}} & ViT-S 3D, parche 8$^3$, 12 bloques & cabezas DINO e iBOT: 384$\rightarrow$2048$\rightarrow$2048$\rightarrow$256$\rightarrow$4096 prototipos & 384 (CLS del profesor) \\
\bottomrule
\end{tabularx}
\caption{Encoders y cabezas. Las cabezas se descartan después del SSL.}
\label{tab:arch}
\end{table}

\begin{figure}[H]
\centering
\begin{tikzpicture}[x=1cm,y=1cm]
% VICReg
\node[font=\small\bfseries, text=vicregc] at (-5.2,0.9) {VICReg};
\node[viewbox, text width=1.6cm] (va) at (-6.2,0) {vista 1};
\node[viewbox, text width=1.6cm] (vb) at (-4.2,0) {vista 2};
\node[netbox=vicregc, text width=3.6cm] (ve) at (-5.2,-1.3) {\textbf{ConvNeXt-T 3D}\\una sola red, mismos pesos};
\node[netbox=vicregc, text width=3.6cm] (vp) at (-5.2,-2.6) {\textbf{Proyector} 2048\\$\Rightarrow z_1,\ z_2$};
\draw[arr] (va) -- (va |- ve.north); \draw[arr] (vb) -- (vb |- ve.north);
\draw[arr] (ve) -- (vp);
% BYOL
\node[font=\small\bfseries, text=byolc] at (0,0.9) {BYOL};
\node[netbox=byolc, text width=2.0cm] (bo) at (-1.35,-1.3) {\textbf{online}\\ResNet-18};
\node[targetbox, text width=2.0cm] (bt) at (1.35,-1.3) {\textbf{target}\\EMA del online};
\node[netbox=byolc, text width=2.0cm] (bp) at (-1.35,-2.6) {proyector +\\\textbf{predictor} $\Rightarrow p$};
\node[targetbox, text width=2.0cm] (btp) at (1.35,-2.6) {proyector\\$\Rightarrow \bar z$ (sin grad.)};
\node[viewbox, text width=1.6cm] (b1) at (-1.35,0) {vista 1};
\node[viewbox, text width=1.6cm] (b2) at (1.35,0) {vista 2};
\draw[arr] (b1) -- (bo); \draw[arrg] (b2) -- (bt);
\draw[arr] (bo) -- (bp); \draw[arrg] (bt) -- (btp);
\draw[emaarr] (bo) -- (bt);
% DINOv2
\node[font=\small\bfseries, text=dinoc] at (5.3,0.9) {DINOv2};
\node[netbox=dinoc, text width=2.0cm] (ds) at (3.95,-1.3) {\textbf{estudiante}\\ViT-S 3D};
\node[targetbox, text width=2.0cm] (dt) at (6.65,-1.3) {\textbf{profesor}\\EMA del estudiante};
\node[netbox=dinoc, text width=2.0cm] (dh) at (3.95,-2.6) {cabezas $\rightarrow$\\4096\\prototipos};
\node[targetbox, text width=2.0cm] (dth) at (6.65,-2.6) {cabezas\\Sinkhorn-Knopp};
\node[viewbox, text width=1.9cm] (d1) at (3.95,0) {vista 1 + máscara};
\node[viewbox, text width=1.6cm] (d2) at (6.65,0) {vista 2};
\draw[arr] (d1) -- (ds); \draw[arrg] (d2) -- (dt);
\draw[arr] (ds) -- (dh); \draw[arrg] (dt) -- (dth);
\draw[emaarr] (ds) -- (dt);
% pérdidas
\node[font=\scriptsize, text=primaryblue, align=center] (lv) at (-5.2,-3.7) {acercar $z_1$ y $z_2$\\sin que colapsen};
\node[font=\scriptsize, text=primaryblue, align=center] (lb) at (0,-3.7) {$p$ predice $\bar z$\\de la otra vista};
\node[font=\scriptsize, text=primaryblue, align=center] (ld) at (5.3,-3.7) {el estudiante copia el\\reparto en prototipos};
\draw[arr] (vp) -- (lv);
\draw[arr] (bp) -- (lb.north west); \draw[arrg] (btp) -- (lb.north east);
\draw[arr] (dh) -- (ld.north west); \draw[arrg] (dth) -- (ld.north east);
\end{tikzpicture}
\caption{Las tres familias. VICReg usa una sola red. BYOL y DINOv2 usan dos copias: la que
aprende (color) y una copia lenta, promedio móvil (EMA) de sus pesos (gris, sin gradiente).
Las flechas naranjas discontinuas son la actualización EMA.}
\label{fig:arch}
\end{figure}

\begin{figure}[!t]
\centering
\includegraphics[width=\textwidth]{architectures.png}
\caption{Los tres modelos campeones, de la entrada a la pérdida. El vector marcado
\emph{probe} es la representación que se evalúa; proyectores, predictor y cabezas solo
existen para entrenar. En BYOL y DINOv2, la caja punteada es la copia lenta (EMA) que
produce los objetivos. Abajo, los hiperparámetros fundamentales de cada uno.}
\label{fig:architectures}
\end{figure}

% ============================================================
\FloatBarrier
\section{Funciones de pérdida}
\label{sec:losses}
% ============================================================

\subsection{VICReg: tres fuerzas}

\begin{equation*}
\mathcal{L}_{\text{VICReg}} = 25\cdot\underbrace{\mathrm{MSE}(z_1,z_2)}_{\text{invarianza}}
+ 25\cdot\underbrace{\big[v(z_1)+v(z_2)\big]}_{\text{varianza}}
+ 1\cdot\underbrace{\big[c(z_1)+c(z_2)\big]}_{\text{covarianza}}
\end{equation*}

La \textbf{invarianza} acerca las dos vistas. La \textbf{varianza} $v$ castiga cada dimensión
cuya desviación estándar en el lote quede bajo 1, así que nadie colapsa a una constante. La
\textbf{covarianza} $c$ castiga dimensiones redundantes. Las dos últimas obligan a separar
sujetos.

\begin{figure}[H]
\centering
\includegraphics[width=\textwidth]{split361_curvas_vicreg.png}
\caption{VICReg: la invarianza cae de 0.23 a 0.06; la varianza es casi toda la pérdida. La
invarianza de validación vale 0 porque el par de validación (cubo y su espejo) es el mismo
input para una red con pesos ImageNet inflados.}
\end{figure}

\subsection{BYOL: predecir a la copia lenta}

\begin{equation*}
\mathcal{L}_{\text{BYOL}} = \tfrac12\Big[\big(2-2\cos(p_1,\bar z_2)\big)+\big(2-2\cos(p_2,\bar z_1)\big)\Big],
\qquad \theta_{\text{target}}\leftarrow\tau\,\theta_{\text{target}}+(1-\tau)\,\theta_{\text{online}}
\end{equation*}

La red online predice, con su predictor, la proyección que la target da a la \emph{otra}
vista. $\tau$ sube de 0.996 a 1 en coseno. Ningún término separa sujetos: lo que evita el
colapso es la asimetría (predictor solo en la online, sin gradiente en la target).

\begin{figure}[H]
\centering
\includegraphics[width=\textwidth]{split361_curvas_byol.png}
\caption{BYOL: el entrenamiento baja a 0.16 y la validación sube desde la época 4 aproximadamente. La
Sección~\ref{sec:byol} muestra que no es memorización.}
\end{figure}

\subsection{DINOv2: repartir en prototipos}

\begin{equation*}
\mathcal{L}_{\text{DINOv2}} = \underbrace{\mathrm{CE}\big(P_t(\text{CLS}_2),\,P_s(\text{CLS}_1)\big)}_{\text{DINO}}
+ 1.0\cdot\underbrace{\mathrm{CE}\big(P_t(\text{parches}),\,P_s(\text{parches enmasc.})\big)}_{\text{iBOT}}
+ 0.1\cdot\underbrace{\mathrm{KoLeo}}_{\text{dispersión}}
\end{equation*}

Cada imagen se reparte en 4096 prototipos (softmax con temperatura 0.1 para el estudiante y
0.01 para el profesor). \textbf{DINO}: el CLS del estudiante reproduce el reparto que el
profesor da a la otra vista. \textbf{iBOT}: los parches tapados reproducen los del profesor.
\textbf{KoLeo}: separa los CLS del lote. Sinkhorn-Knopp reparte los prototipos de forma pareja
para que no se concentren en uno.

\begin{figure}[H]
\centering
\includegraphics[width=\textwidth]{split361_curvas_dinov2.png}
\caption{DINOv2: DINO baja de 8.4 a 6.75 (bajo $\ln K = 8.3$), iBOT de 4.2 a 3.4; a diferencia
del informe v2, la validación ahora también baja. La entropía del profesor se mantiene en
0.67--0.76\,$\ln K$: sin colapso.}
\end{figure}

\subsection{Lo común a los tres}

AdamW con tasa $10^{-4}$, 6 épocas de \emph{warmup} y decaimiento coseno hasta $10^{-6}$; lote
de 16; 30 épocas; bf16; 3 semillas; se evalúa el estado final, sin \emph{early stopping}.
DINOv2 usa además \emph{weight decay} 0.04$\rightarrow$0.2 y recorte de gradiente.

% ============================================================
\section{Cómo evaluamos}
\label{sec:eval}
% ============================================================

\begin{figure}[H]
\centering
\begin{tikzpicture}[x=1cm,y=1cm]
\node[evalbox] (enc) at (0,0) {\textbf{Encoder congelado}\\cabezas descartadas};
\node[evalbox] (emb) at (4.2,0) {\textbf{Un vector por sujeto}\\promedio de sus sesiones,\\cubo limpio};
\node[evalbox] (lr) at (8.4,0) {\textbf{Regresión logística}\\estandarizada, clases balanceadas};
\node[evalbox, fill=white] (cv) at (2.1,-2.3) {\textbf{CV 5$\times$5}\\329 sujetos AD/HC\\(266 train + 63 val)};
\node[evalbox, fill=white] (te) at (6.3,-2.3) {\textbf{Test}\\ajuste en 266 de train,\\$C$ elegido en val, una vez\\sobre 84 sujetos};
\node[evalbox, fill=white, draw=accentorange] (col) at (10.5,-2.3) {\textbf{Monitor de colapso}\\cada 2 épocas, en 147\\sujetos no vistos};
\draw[arr] (enc) -- (emb); \draw[arr] (emb) -- (lr);
\draw[arr] (lr.south) -- ++(0,-0.5) -| (cv.north);
\draw[arr] (lr.south) -- ++(0,-0.5) -| (te.north);
\draw[emaarr] (enc.south) -- ++(0,-0.25) -| (col.north);
\end{tikzpicture}
\caption{Del encoder entrenado a los números del informe.}
\label{fig:eval}
\end{figure}

\begin{itemize}
\item \textbf{Métricas:} AUROC (la principal), F1-macro y AUPR.
\item \textbf{Comparación justa:} cada modelo se compara con \emph{el mismo encoder y la misma
  entrada sin SSL} (pesos ImageNet o DeiT inflados). La diferencia lleva un intervalo de
  confianza del 95\,\% por \emph{bootstrap} pareado sobre sujetos.
\item \textbf{Colapso:} rango efectivo, varianza entre sujetos sobre varianza por augmentation y
  top-1 (si una vista encuentra a su gemela entre 147; azar 0.7\,\%).
\end{itemize}

\begin{warnbox}[CV y test no miden lo mismo]
266 de los 329 sujetos del CV estuvieron en el entrenamiento SSL (sin etiquetas). El test son
sujetos que el encoder nunca vio, pero con 84 sujetos sus intervalos son anchos (unos
$\pm0.05$ a $\pm0.07$).
\end{warnbox}

% ============================================================
\section{Resultados}
\label{sec:results}
% ============================================================

\begin{figure}[H]
\centering
\includegraphics[width=\textwidth]{split361_metricas.png}
\caption{Las tres métricas en CV y test. Color: modelo con SSL; gris: mismo encoder sin SSL;
puntos: semillas.}
\label{fig:metrics}
\end{figure}

\begin{figure}[H]
\centering
\includegraphics[width=0.9\textwidth]{delta.png}
\caption{Lo que aporta el SSL: diferencia de AUROC con el mismo encoder sin entrenar, con IC
del 95\,\%. Un intervalo que no toca el cero es una mejora significativa.}
\label{fig:delta}
\end{figure}

\begin{ideabox}[Cómo leerlo]
\begin{itemize}[leftmargin=12pt]
\item \textbf{DINOv2} es el único que mejora en CV \emph{y} en test, pero porque su ViT sin
  entrenar es el peor punto de partida (0.778 CV, 0.727 test). Con SSL alcanza a las CNN, no
  las supera.
\item \textbf{VICReg} mejora en CV; en test, $+0.029$ queda dentro del ruido.
\item \textbf{BYOL} queda igual que su ResNet ImageNet, que ya era el mejor encoder sin SSL
  (0.860 CV).
\end{itemize}
\end{ideabox}

% ============================================================
\section{¿Colapsan?}
\label{sec:collapse}
% ============================================================

No. En los 147 sujetos no vistos, la diferencia entre sujetos crece a lo largo del
entrenamiento y termina muy por encima de la variación por augmentation.

\begin{figure}[H]
\centering
\includegraphics[width=0.9\textwidth]{collapse.png}
\caption{Encoder online, media de 3 semillas. Izquierda: sobre 1 (línea punteada), el sujeto
pesa más que la augmentation. Derecha: fracción de vistas que encuentran a su gemela.}
\label{fig:collapse}
\end{figure}

\begin{table}[H]
\centering
\small
\begin{tabular}{lrrl}
\toprule
Modelo & Var sujeto / Var aug. & Rango efectivo & Top-1 de vistas \\
\midrule
VICReg & 23.2 & 7.2 & 0.75 \\
BYOL & 34.2 & 10.5 & 0.96 \\
DINOv2 & 10.5 & 22.7 & 1.00 \\
\bottomrule
\end{tabular}
\caption{Valores finales (época 30). En el informe v2, BYOL tenía cocientes de $4\times10^{-5}$
y rangos de 1--1.2.}
\end{table}

\begin{warnbox}[Reconocer al sujeto no es reconocer la enfermedad]
Los tres aprenden muy bien \emph{quién} es cada sujeto (DINOv2 reconoce el 100\,\% de las
vistas), pero el probe mejora poco. La identidad del sujeto no basta para separar AD de HC.
\end{warnbox}

% ============================================================
\section{El espacio latente}
\label{sec:latent}
% ============================================================

Los embeddings salen directo de los modelos campeones de este informe (semilla 1, estado final),
sin reentrenar nada: un vector por sujeto (promedio de sus sesiones, cubo limpio) para los 413
sujetos AD/HC. El probe de esta sección es una regresión logística evaluada \emph{fuera de
muestra} con 5 pliegues sobre los 413 sujetos, así que sus AUROC difieren algo de los de la
Sección~\ref{sec:results}.

\subsection{AD y HC en el espacio latente}

\begin{figure}[H]
\centering
\includegraphics[width=0.95\textwidth]{latent.png}
\caption{Cada punto es un sujeto. Fila 1: las dos direcciones de mayor varianza (PCA). Fila 2:
t-SNE, que conserva vecindarios. Fila 3: puntaje del probe para cada sujeto, calculado sin que
el probe lo haya visto.}
\label{fig:latent}
\end{figure}

\begin{itemize}
\item \textbf{AD y HC no forman grupos separados}, sino un gradiente: en t-SNE los vecinos
  tienden a compartir diagnóstico, pero hay mucha mezcla.
\item \textbf{La dirección principal ya ordena el diagnóstico}: en VICReg y BYOL la primera
  componente (32--34\,\% de la varianza) correlaciona $-0.56$ con AD; en DINOv2, $+0.51$
  (el signo de una componente es arbitrario).
\item \textbf{El probe separa con solapamiento}: AUROC fuera de muestra de 0.86 (VICReg), 0.87
  (BYOL) y 0.80 (DINOv2).
\end{itemize}

\subsection{¿Qué codifican los embeddings?}

\begin{figure}[H]
\centering
\includegraphics[width=\textwidth]{latent_corr.png}
\caption{Correlación de cada variable con las 5 primeras componentes del embedding y con su
mejor predicción lineal desde el vector completo. El ruido se mide en la sustancia blanca de la
resonancia completa (dispersión de alta frecuencia en un tejido homogéneo). Sexo no está
disponible en OASIS3.}
\label{fig:latentcorr}
\end{figure}

\begin{figure}[p]
\centering
\includegraphics[width=0.9\textwidth]{latent_tsne_vars.png}
\caption{El mismo mapa t-SNE de cada modelo, coloreado por cada variable (color = percentil
entre los 413 sujetos). Una variable que el embedding codifica aparece como un gradiente suave
sobre el mapa; una que no codifica, como sal y pimienta (educación). $r$: correlación entre la
variable y su predicción lineal desde el embedding completo, fuera de muestra (la última
columna de la Figura~\ref{fig:latentcorr}). Con marco de color, las variables que cada modelo
codifica al menos tan bien como el diagnóstico: todas menos la educación.}
\label{fig:tsnevars}
\end{figure}

\begin{itemize}
\item \textbf{La primera componente es un eje de atrofia}: sigue al volumen del hipocampo
  ($|\rho|$ 0.71--0.84), al de los ventrículos (0.70--0.80) y a la edad (0.53--0.60).
\item \textbf{El vector completo codifica casi perfecto la anatomía}: volumen del hipocampo
  (0.87--0.91), ventrículos (0.83--0.86) y volumen cerebral (0.81--0.83). La educación no
  aparece ($\leq$0.10).
\item \textbf{También codifica la intensidad y el ruido} (0.76--0.96): el encoder ve toda la
  imagen, no solo lo útil para el diagnóstico.
\end{itemize}

\begin{warnbox}[La edad pesa tanto como el diagnóstico]
El embedding predice la edad (0.57--0.67) tan bien como el diagnóstico (0.51--0.65). Como los AD
son 7 años mayores (Sección~\ref{sec:data}), parte del eje AD es un eje de envejecimiento.
\end{warnbox}

\subsection{¿Cómo garantizamos que no aprende solo el ruido?}

No hay una prueba única. Lo que sí se puede hacer es buscar la huella que dejaría un modelo que
aprendió ruido y comprobar que no aparece. Hay tres huellas fáciles de medir.

\begin{figure}[H]
\centering
\includegraphics[width=\textwidth]{noise_checks.png}
\caption{Tres pruebas sobre los embeddings de los campeones. 1: probe ajustado en train+val con
cubos limpios y evaluado en los 84 sujetos de test con ruido gaussiano añadido. 2: para cada
sesión, ¿su vecino más cercano entre 131 sesiones es otra sesión del mismo sujeto? 3:
correlación parcial de Spearman del puntaje AD (marca negra: la del diagnóstico real).}
\label{fig:noise}
\end{figure}

\begin{enumerate}[leftmargin=16pt]
\item \textbf{Si dependiera del ruido, más ruido cambiaría la respuesta.} Con ruido añadido
  hasta $\sigma=0.3$, los modelos con SSL mantienen el AUROC de test (VICReg 0.83$\rightarrow$0.86,
  BYOL 0.87$\rightarrow$0.86, DINOv2 0.85$\rightarrow$0.80). Los mismos encoders sin SSL caen
  (BYOL 0.82$\rightarrow$0.59, DINOv2 0.76$\rightarrow$0.53): el SSL los hizo más robustos al
  ruido.
\item \textbf{Si codificara el ruido de cada imagen, no reconocería al sujeto en otra
  sesión.} Dos resonancias de la misma persona en días distintos tienen ruido distinto y la
  misma anatomía. El vecino más cercano es otra sesión del mismo sujeto el 23\,\% (VICReg), 47\,\%
  (BYOL) y 79\,\% (DINOv2) de las veces, frente a un azar de 1\,\%. VICReg y BYOL quedan bajo sus
  encoders sin SSL (67\,\% y 82\,\%): su invarianza descarta parte del detalle fino.
\item \textbf{Si siguiera al ruido, el eje AD no seguiría a la anatomía.} Descontando edad y
  ruido, el puntaje AD correlaciona con el volumen del hipocampo ($|\rho|$ 0.62--0.77), más que el
  propio diagnóstico (0.54). En DINOv2, borrar el canal del hipocampo baja el AUROC de test de
  0.85 a 0.77; borrar el ventrículo o la corteza posteromedial apenas lo mueve (0.84--0.85).
\end{enumerate}

\begin{ideabox}[Respuesta corta]
Las tres pruebas apuntan a que la decisión se apoya en la \textbf{anatomía} (sobre todo en el
hipocampo) y no en el ruido. Queda una señal residual: descontando anatomía y edad, el puntaje AD
todavía correlaciona 0.18--0.28 con el ruido de la sustancia blanca, más que el diagnóstico real
(0.10). Puede ser textura de lesiones de sustancia blanca o algo de la adquisición. Para
\emph{garantizarlo} faltan dos pruebas: validar en otra cohorte u otro escáner (ADNI ya está en
el LMDB) y controlar por escáner, un dato que hoy no está en el CSV.
\end{ideabox}

% ============================================================
\section{BYOL: ¿por qué sube su validación?}
\label{sec:byol}
% ============================================================

La pérdida de validación de BYOL sube desde la época 4 mientras la de entrenamiento baja.
Parecía sobreajuste. Se probó con lotes más grandes y con GroupNorm en las cabezas (30 épocas,
3 semillas), y se midió el objetivo con las mismas vistas en sujetos de train y no vistos.

\begin{figure}[H]
\centering
\includegraphics[width=\textwidth]{byol_diag_objetivo.png}
\caption{Paneles 3 y 4: con la red en modo evaluación, el objetivo es idéntico en sujetos de
train y no vistos (brecha bajo 0.015). No hay memorización.}
\label{fig:byoldiag}
\end{figure}

\begin{table}[H]
\centering
\small
\begin{tabular}{lrrrr}
\toprule
Configuración & AUROC CV & AUROC test & Objetivo (modo eval.) & Subida de val. \\
\midrule
Actual (lote 16, BN) & 0.851 & 0.831 & 1.39 & 0.42 \\
Lote 32 & 0.847 & 0.833 & 1.18 & 0.34 \\
Lote 64 & 0.848 & 0.829 & 1.12 & 0.17 \\
GroupNorm en cabezas & 0.853 & 0.821 & 0.35 & 0.09 \\
\bottomrule
\end{tabular}
\caption{Ninguna variante cambia el probe: diferencias de AUROC entre $-0.010$ y $+0.003$,
ninguna significativa. Objetivo y subida: pérdida $2-2\cos$ al final del entrenamiento.}
\end{table}

\begin{byolbox}[Diagnóstico]
No es memorización: es la \textbf{batch norm}. En entrenamiento las cabezas normalizan con las
estadísticas del lote (pérdida 0.16); en evaluación, con promedios globales, y ahí el objetivo
vale 1.4 \emph{incluso en sujetos de train}. Lotes más grandes achican la brecha y GroupNorm casi
la elimina, pero el encoder que usa el probe rinde igual. La curva que sube es inofensiva.
\end{byolbox}

% ============================================================
\section{Costo y velocidad}
\label{sec:cost}
% ============================================================

\begin{table}[H]
\centering
\small
\begin{tabular}{lrrrr}
\toprule
Modelo & volúmenes/s & s por época & VRAM pico & min por corrida \\
\midrule
BYOL & 790 & 1.1 & 10.0 GB & 1.8 \\
VICReg & 469 & 1.7 & 1.1 GB & 2.0 \\
DINOv2 & 27 & 27.1 & 1.3 GB & 16.6 \\
\bottomrule
\end{tabular}
\caption{RTX 5090 (RunPod, 0.69 USD/h), 9 corridas en paralelo con CUDA MPS: 20.2 minutos de
reloj y 0.33 USD.}
\end{table}

DINOv2 es 16 a 25 veces más lento por época porque la deformación elástica de TorchIO se calcula en la CPU,
no por el ViT: la GPU estuvo ocupada solo el 13\,\% del tiempo.

% ============================================================
\section{Entrenamiento más largo: 100 épocas}
\label{sec:e100}
% ============================================================

Repetimos el entrenamiento del informe con \textbf{100 épocas en vez de 30}: mismo split, mismas
entradas, mismas augmentations, 3 semillas, misma evaluación. Tres cambios:

\begin{itemize}
\item \textbf{BYOL usa la variante elegida en la Sección~\ref{sec:byol}}: GroupNorm en proyector y
  predictor, dropout 0.3 en el proyector y weight decay 0.01 (\texttt{gn\_dp0.3\_dq0\_wd0.01}).
\item \textbf{VICReg valida con vistas aumentadas.} Antes la pérdida de validación comparaba el cubo
  con su espejo L/R; como el entrenamiento ya incluye ese espejo, la invarianza de validación
  quedaba en cero y no informaba nada. Ahora usa las augmentations de entrenamiento, fijas por
  sesión (la misma semilla en cada época), así que la curva es comparable entre épocas. La
  evaluación (probe) sigue con el cubo limpio.
\item \textbf{Dos referencias nuevas para cada encoder} (Tabla~\ref{tab:refs100}).
\end{itemize}

\begin{table}[H]
\centering
\small
\begin{tabularx}{\textwidth}{lX}
\toprule
Barra & Qué es \\
\midrule
\textbf{Supervisado completo} & El mismo encoder (mismos pesos iniciales, entrada y augmentation)
  más una capa lineal, entrenados de punta a punta con BCE balanceada en los 266 sujetos AD/HC
  de train (AdamW, lr $10^{-4}$, warmup y coseno, lote 16). Se queda la época con mejor AUROC de
  validación (máx.\ 100, paciencia 20). CV: 5 folds sobre los 329 sujetos, cada uno entrenado
  hasta esa época. \\
\textbf{SSL} & El encoder tras 100 épocas de SSL, congelado, con el probe de la
  Sección~\ref{sec:eval}. \\
\textbf{Sin SSL} & El mismo encoder \emph{antes} del SSL (pesos ImageNet o DeiT inflados, el
  punto de partida del SSL), congelado, con el mismo probe. Sus pesos no dependen de la semilla
  (solo el \emph{mask token} del ViT, que no se usa al evaluar), así que una corrida vale para
  las tres. \\
\bottomrule
\end{tabularx}
\caption{Las tres barras de la Figura~\ref{fig:metrics100}, en su orden. Como referencia extra
también evaluamos el mismo encoder con pesos aleatorios
(\nolinkurl{split361_100/REPORT.md}).}
\label{tab:refs100}
\end{table}

\begin{warnbox}[El supervisado necesita calibración]
Cada modelo supervisado termina con su propio sesgo: detenidos en una época temprana, la mediana
de su probabilidad de AD va de 0.0 a 0.98 según el fold. Juntar las predicciones de los 5 folds
bajaba el AUROC de CV de unos 0.84 (dentro de cada fold) a 0.74, y el umbral 0.5 del F1 dejaba de
tener sentido. Por eso cada modelo supervisado recibe la misma calibración que el probe: una
regresión logística balanceada sobre sus propios logits en sus sujetos de entrenamiento (sin
etiquetas de validación ni de test). No cambia el AUROC de test; corrige el CV y el F1.
\end{warnbox}

\subsection{Las tres arquitecturas, rama por rama}

Cada figura separa las dos ramas en bandas de distinto color. A la izquierda, el cubo limpio de un
sujeto de validación y las dos vistas que su receta produce (cortes centrales del volumen de
32$^3$). En VICReg las dos ramas son la misma red; en BYOL y DINOv2 la rama de abajo (punteada) es
una copia lenta (EMA) que no recibe gradiente (\emph{sg}).

\begin{figure}[H]
\centering
\includegraphics[width=\textwidth]{arch_vicreg.png}
\caption{VICReg: dos vistas \texttt{geo} por la misma ConvNeXt-T 3D; la pérdida acerca $z_1$ y
$z_2$ y mantiene la varianza y la decorrelación de las dimensiones.}
\label{fig:archvicreg}
\end{figure}

\begin{figure}[H]
\centering
\includegraphics[width=\textwidth]{arch_byol.png}
\caption{BYOL con la variante elegida: cabezas con GroupNorm, dropout 0.3 en el proyector y
weight decay 0.01. Solo la rama online tiene predictor y recibe gradiente.}
\label{fig:archbyol}
\end{figure}

\begin{figure}[H]
\centering
\includegraphics[width=\textwidth]{arch_dinov2.png}
\caption{DINOv2: el alumno ve la vista con un bloque de parches ocultos (color) y el profesor la
vista completa; el probe usa el CLS del profesor.}
\label{fig:archdinov2}
\end{figure}

\subsection{Métricas}

\begin{figure}[H]
\centering
\includegraphics[width=\textwidth]{split361_100_metricas.png}
\caption{Las tres métricas tras 100 épocas. Rayado: supervisado completo; color: SSL; gris: el
mismo encoder antes del SSL. Puntos: semillas.}
\label{fig:metrics100}
\end{figure}

\begin{figure}[H]
\centering
\includegraphics[width=0.85\textwidth]{delta_100.png}
\caption{Diferencia de AUROC del SSL frente al mismo encoder sin SSL (círculo lleno) y frente al
supervisado (cuadrado vacío), con IC del 95\,\% por bootstrap pareado.}
\label{fig:delta100}
\end{figure}

\begin{table}[H]
\centering
\small
\begin{tabular}{lrrrrrr}
\toprule
 & \multicolumn{4}{c}{AUROC CV / test} & \multicolumn{2}{c}{$\Delta$ AUROC del SSL 100 (CV / test)} \\
\cmidrule(lr){2-5}\cmidrule(lr){6-7}
Modelo & Supervisado & SSL 100 & SSL 30 & Sin SSL & vs sin SSL & vs supervisado \\
\midrule
\textcolor{vicregc}{\textbf{VICReg}} & .864 / .849 & .877 / .834 & .875 / .835 & .848 / .806 &
  $+.028^{*}$ / $+.028$ & $+.016$ / $-.015$ \\
\textcolor{byolc}{\textbf{BYOL}} & .827 / .850 & .846 / .838 & .854 / .843 & .860 / .808 &
  $-.015$ / $+.030$ & $+.026^{*}$ / $-.012$ \\
\textcolor{dinoc}{\textbf{DINOv2}} & .833 / .825 & .832 / .833 & .839 / .831 & .778 / .727 &
  $+.047^{*}$ / $+.106^{*}$ & $+.002$ / $+.009$ \\
\bottomrule
\end{tabular}
\caption{Media de 3 semillas. $^{*}$ = IC del 95\,\% enteramente sobre cero. De 30 a 100 épocas
ninguna diferencia supera $\pm0.008$ ni es significativa. Detalle en
\nolinkurl{campaign_ssl3d_runpod/split361_100/REPORT.md}.}
\label{tab:summary100}
\end{table}

\begin{ideabox}[Cómo leerlo]
\begin{itemize}[leftmargin=12pt]
\item \textbf{Entrenar más no ayuda}: con 100 épocas los tres quedan donde estaban a las 30.
\item \textbf{Frente al mismo encoder sin SSL}, DINOv2 gana mucho ($+0.047$ CV, $+0.106$ test,
  ambos significativos) porque su ViT de partida es el peor (0.778 / 0.727). VICReg gana en CV
  ($+0.028$, significativo) y $+0.028$ en test (dentro del ruido). BYOL no supera a su ResNet
  ImageNet en CV ($-0.015$); en test, $+0.030$ no es significativo.
\item \textbf{Frente al supervisado completo}, el SSL empata en test (de $-0.015$ a $+0.009$, nada
  significativo) y empata o gana en CV. Con 266 sujetos etiquetados, el SSL con un probe lineal
  rinde lo mismo que entrenar la red entera con las etiquetas.
\item \textbf{El supervisado tampoco se separa mucho del encoder de partida}, salvo en el ViT
  ($+0.045$ CV, $+0.097$ test, significativos). En las dos CNN gana $+0.043$ en test sin
  significancia, y en la ResNet de BYOL queda incluso $0.040$ bajo ella en CV (significativo).
\end{itemize}
\end{ideabox}

\begin{warnbox}[El CV del supervisado no es el mismo CV]
El probe repite 5 veces una CV de 5 folds; el supervisado hace una sola (5 entrenamientos por
semilla) y cada fold entrena hasta la época que eligió la validación. El test (mismos 84 sujetos,
mismo modelo final) es la comparación más limpia.
\end{warnbox}

\subsection{Curvas de entrenamiento}

\begin{figure}[H]
\centering
\includegraphics[width=\textwidth]{split361_100_curvas_vicreg.png}
\includegraphics[width=0.8\textwidth]{split361_100_curvas_byol.png}
\includegraphics[width=\textwidth]{split361_100_curvas_dinov2.png}
\caption{Pérdidas de los tres modelos en 100 épocas (continua: train; discontinua: validación).}
\label{fig:curves100}
\end{figure}

\begin{itemize}
\item \textbf{VICReg}: la invarianza de validación ya no es cero. Sube a 0.05 en la primera época y
  baja hasta 0.007: el encoder se vuelve invariante a las augmentations también en sujetos que no
  entrena. Queda bajo la de train (0.04) porque en validación las dos vistas se normalizan con las
  mismas estadísticas de batch norm, y en train cada vista con las de su lote.
\item \textbf{BYOL}: la validación (cubo vs espejo) baja hasta la época 13--17 (0.13--0.19) y
  luego sube y se estabiliza en 0.24--0.53 según la semilla. La pérdida de train sigue bajando
  hasta 0.02.
\item \textbf{DINOv2}: las tres pérdidas siguen bajando hasta el final; la de validación se separa
  de la de train desde la época 20 y la entropía del profesor se estabiliza en 0.65 de su máximo.
\end{itemize}

\subsection{Colapso}

\begin{figure}[H]
\centering
\includegraphics[width=0.9\textwidth]{collapse_100.png}
\caption{Igual que la Figura~\ref{fig:collapse}, en 100 épocas (punteada: la corrida de 30).}
\label{fig:collapse100}
\end{figure}

\begin{table}[H]
\centering
\small
\begin{tabular}{lrrrrrr}
\toprule
 & \multicolumn{2}{c}{Var sujeto / Var aug.} & \multicolumn{2}{c}{Rango efectivo} &
 \multicolumn{2}{c}{Top-1 de vistas} \\
Modelo & 30 ép. & 100 ép. & 30 ép. & 100 ép. & 30 ép. & 100 ép. \\
\midrule
VICReg & 23.2 & 47.1 & 7.2 & 9.5 & 0.75 & 0.95 \\
BYOL$^{\dagger}$ & 34.2 & 56.7 & 10.5 & 13.2 & 0.96 & 0.99 \\
DINOv2 & 10.5 & 17.4 & 22.7 & 25.6 & 1.00 & 1.00 \\
\bottomrule
\end{tabular}
\caption{Valores finales en los 147 sujetos no vistos. $^{\dagger}$BYOL a 30 épocas es la
variante anterior (\texttt{wide\_dproj0.1}).}
\end{table}

Ninguno colapsa: con más épocas la diferencia entre sujetos sigue creciendo en VICReg y BYOL, y en
DINOv2 se estabiliza desde la época 30. Todos reconocen casi el 100\,\% de las vistas.

\subsection{BYOL en 100 épocas}

\begin{figure}[H]
\centering
\includegraphics[width=\textwidth]{byol_objetivo_100.png}
\caption{BYOL elegido, 100 épocas (naranja) frente a 30 (gris). Derecha: el objetivo con la red en
modo evaluación y las mismas vistas, en sujetos de train (continua) y no vistos (discontinua).}
\label{fig:byol100}
\end{figure}

\begin{byolbox}[Con más épocas aparece una brecha pequeña]
A 30 épocas el objetivo en modo evaluación era el mismo en sujetos de train y no vistos (brecha
$\leq0.015$). A 100 épocas los no vistos quedan 0.03--0.07 por encima desde la época 60: ahora
sí hay un poco de memorización del objetivo, pero no se traduce en el probe.
\end{byolbox}

\subsection{El supervisado}

\begin{figure}[H]
\centering
\includegraphics[width=0.9\textwidth]{supervised_100.png}
\caption{Entrenamiento supervisado, una línea por semilla. Punto: la época elegida en
validación.}
\label{fig:sup100}
\end{figure}

\begin{itemize}
\item \textbf{Aprende rápido y sobreajusta pronto}: la ResNet-18 (BYOL) llega a su mejor
  validación en las épocas 6--13 y su BCE de entrenamiento cae hasta 0.05; la ConvNeXt (VICReg)
  en las 12--49; el ViT (DINOv2) en las 18--51, con una semilla estancada en BCE 0.69 durante
  las primeras 45 épocas.
\item \textbf{En test es el mejor número del informe} (0.849 ConvNeXt, 0.850 ResNet, 0.825 ViT),
  pero dentro del ruido del SSL (Tabla~\ref{tab:summary100}).
\item \textbf{Frente al encoder sin SSL} (el de partida, congelado) solo el ViT mejora de forma
  significativa ($+0.045$ CV, $+0.097$ test); la ConvNeXt gana $+0.043$ en test sin
  significancia y la ResNet queda $0.040$ por debajo en CV.
\end{itemize}

\subsection{Costo}

RTX 5090 en RunPod. SSL y encoders aleatorios: 1.83 horas, 1.27 USD (DINOv2 tomó 102 minutos
por semilla, 3 en paralelo; VICReg y BYOL, 8--10 minutos). Supervisado: 0.82 USD en un host con
red lenta, más 0.07 USD en tres pods descartados porque CUDA no arrancaba en su contenedor.
Total: 2.16 USD.

\subsection{El espacio latente tras 100 épocas}

Mismo análisis de la Sección~\ref{sec:latent}, ahora con los modelos de 100 épocas (semilla 1,
estado final). ``Sin SSL'' es, como en las métricas, el encoder \emph{de partida} (pesos ImageNet
o DeiT inflados), el punto desde el que arranca el SSL.

\begin{figure}[H]
\centering
\includegraphics[width=0.95\textwidth]{latent_100.png}
\caption{Como la Figura~\ref{fig:latent}, con los modelos de 100 épocas.}
\label{fig:latent100}
\end{figure}

\begin{figure}[H]
\centering
\includegraphics[width=\textwidth]{latent_corr_100.png}
\caption{Como la Figura~\ref{fig:latentcorr}, con los modelos de 100 épocas.}
\label{fig:latentcorr100}
\end{figure}

\begin{figure}[p]
\centering
\includegraphics[width=0.9\textwidth]{latent_tsne_vars_100.png}
\caption{Como la Figura~\ref{fig:tsnevars}, con los modelos de 100 épocas.}
\label{fig:tsnevars100}
\end{figure}

\begin{figure}[H]
\centering
\includegraphics[width=\textwidth]{noise_checks_100.png}
\caption{Como la Figura~\ref{fig:noise}, con los modelos de 100 épocas.}
\label{fig:noise100}
\end{figure}

\begin{table}[H]
\centering
\small
\begin{tabular}{lrrrrrr}
\toprule
 & \multicolumn{2}{c}{VICReg} & \multicolumn{2}{c}{BYOL$^{\dagger}$} & \multicolumn{2}{c}{DINOv2} \\
 & 30 & 100 & 30 & 100 & 30 & 100 \\
\midrule
AUROC del probe fuera de muestra & 0.86 & 0.86 & 0.87 & 0.86 & 0.80 & 0.83 \\
Varianza de la 1.ª componente & 0.32 & 0.22 & 0.34 & 0.22 & 0.15 & 0.12 \\
$|\rho|$ 1.ª componente -- AD & 0.56 & 0.49 & 0.56 & 0.59 & 0.51 & 0.56 \\
$|\rho|$ 1.ª componente -- vol.\ hipocampo & 0.71 & 0.67 & 0.84 & 0.84 & 0.72 & 0.76 \\
AUROC test con ruido $\sigma=0.3$ & 0.86 & 0.84 & 0.86 & 0.84 & 0.80 & 0.84 \\
Otra sesión del mismo sujeto (top-1) & 0.23 & 0.23 & 0.47 & 0.65 & 0.79 & 0.74 \\
Puntaje AD -- vol.\ hipocampo (parcial) & 0.77 & 0.76 & 0.76 & 0.79 & 0.62 & 0.77 \\
Puntaje AD -- ruido (parcial) & 0.24 & 0.23 & 0.18 & 0.23 & 0.28 & 0.39 \\
\bottomrule
\end{tabular}
\caption{Lo que cambia en el espacio latente de 30 a 100 épocas (semilla 1). $^{\dagger}$BYOL a 30
épocas es la variante anterior.}
\label{tab:latent100}
\end{table}

\begin{itemize}
\item \textbf{Casi nada cambia en lo que se puede separar}: el probe fuera de muestra queda en
  0.86 (VICReg y BYOL) y DINOv2 sube de 0.80 a 0.83. La primera componente sigue siendo un eje de
  atrofia ($|\rho|$ 0.67--0.84 con el volumen del hipocampo) y explica menos varianza (0.12--0.22):
  el embedding reparte la información en más direcciones, como ya indicaba el rango efectivo.
\item \textbf{La robustez al ruido se mantiene}: con $\sigma=0.3$ los tres quedan en 0.84 de AUROC
  de test, mientras los encoders de partida caen hasta 0.53--0.59.
\item \textbf{El eje AD sigue apoyado en la anatomía}: descontando edad y ruido, el puntaje AD
  correlaciona 0.76--0.79 con el volumen del hipocampo (el diagnóstico real, 0.54); DINOv2 sube
  de 0.62 a 0.77. Pero en DINOv2 la señal residual de ruido también crece, de 0.28 a 0.39
  (VICReg y BYOL quedan en 0.23; el diagnóstico real, 0.10).
\item \textbf{BYOL reconoce mejor al sujeto en otra sesión} (47\,\% $\rightarrow$ 65\,\%), más
  cerca de su encoder de partida (82\,\%): la variante con GroupNorm descarta menos detalle fino.
\end{itemize}

% ============================================================
\section{Conclusiones}
% ============================================================

\begin{enumerate}[leftmargin=16pt]
\item Con el split 361/63/84 los tres métodos son estables: ninguno colapsa y las semillas
  coinciden.
\item En test los tres quedan en 0.83--0.84 de AUROC: no hay un ganador claro.
\item El SSL aporta poco sobre un buen punto de partida: $+0.026$ en VICReg, nada en BYOL.
  DINOv2 gana mucho, pero porque arranca desde abajo.
\item Los encoders aprenden la identidad del sujeto mucho mejor que la enfermedad.
\item La validación creciente de BYOL es un efecto de la batch norm, no sobreajuste. A 100 épocas
  aparece además una brecha pequeña entre sujetos de train y no vistos (0.03--0.07).
\item Entrenar 100 épocas en vez de 30 no cambia ninguna métrica: los modelos ya estaban
  saturados.
\item Con 266 sujetos etiquetados, el SSL con probe lineal rinde lo mismo que el entrenamiento
  supervisado completo del mismo encoder. Frente al mismo encoder sin SSL, DINOv2 gana mucho
  (parte del ViT más débil), VICReg algo (solo en CV) y BYOL nada.
\end{enumerate}

\appendix
\section{Reproducir este informe}

\begin{itemize}
\item Split: \texttt{python} \nolinkurl{scripts/make_split_361_63_84.py}
  $\rightarrow$ \nolinkurl{campaign_ssl3d_runpod/split_361_63_84/}.
\item Entrenamiento: \nolinkurl{run_all.py} con \texttt{-{}-set}
  \nolinkurl{data.manifest_path=campaign_ssl3d_runpod/split_361_63_84/manifest.csv};
  resultados en \nolinkurl{campaign_ssl3d_runpod/split361_30/}.
\item Figuras: \texttt{python scripts/build\_report\_split361.py}; luego, desde \texttt{docs/},
  dos veces \texttt{pdflatex informe\_split361.tex}.
\item 100 épocas (Sección~\ref{sec:e100}): \nolinkurl{run_all.py} con \texttt{-{}-epochs 100
  -{}-identity -{}-collapse-every 2} y las especificaciones
  \nolinkurl{vicreg:geo:32:campaign_augval}, \nolinkurl{byol:noise:32:gn_dp0.3_dq0_wd0.01} y
  \nolinkurl{dinov2:elastic:multi3}; encoders aleatorios con \nolinkurl{train.py --init-only}
  y \nolinkurl{model.backbone.pretrained=false}; supervisado con
  \nolinkurl{scripts/train_supervised.py}. Resultados en
  \nolinkurl{campaign_ssl3d_runpod/split361_100/} (\texttt{REPORT.md}).
\item Figuras de 100 épocas: \nolinkurl{scripts/report_split361_100.py} y
  \nolinkurl{scripts/analyze_latent_split361.py} con \texttt{-{}-runs split361\_100 -{}-out
  latent\_split361\_100 -{}-suffix \_100} y los nombres de VICReg y BYOL en \texttt{-{}-run}.
\end{itemize}

\end{document}
